\documentclass[11pt,a4paper]{article}
\usepackage[margin=2.4cm]{geometry}
\usepackage{ctex}
\usepackage{amsmath}
\usepackage{booktabs}
\usepackage{array}
\usepackage{enumitem}
\usepackage{xcolor}
\usepackage[colorlinks=true,linkcolor=black,urlcolor=teal]{hyperref}
\usepackage{titlesec}
\usepackage{tikz}
\usetikzlibrary{positioning,arrows,shapes.geometric,calc,fit,backgrounds}
\tikzset{
  box/.style={draw=ink!55, fill=white, rounded corners=1.5pt, align=center,
              inner sep=5pt, font=\small, minimum height=9mm},
  dark/.style={draw=ink, fill=ink, text=white, rounded corners=1.5pt, align=center,
              inner sep=5pt, font=\small\bfseries, minimum height=9mm},
  acc/.style={draw=acc, fill=acc!8, rounded corners=1.5pt, align=center,
              inner sep=5pt, font=\small, minimum height=9mm},
  warnb/.style={draw=warn, fill=warn!8, rounded corners=1.5pt, align=center,
              inner sep=5pt, font=\small, minimum height=9mm},
  fl/.style={->, >=latex, draw=ink!60, thick},
  cut/.style={draw=dead, very thick, dashed},
  lbl/.style={font=\footnotesize\itshape, text=ink!70},
}

\definecolor{ink}{HTML}{141A22}
\definecolor{acc}{HTML}{0E6E78}
\definecolor{warn}{HTML}{9A5B00}
\definecolor{dead}{HTML}{A8231C}

\titleformat{\section}{\large\bfseries\color{ink}}{\thesection}{0.6em}{}
\titleformat{\subsection}{\normalsize\bfseries\color{ink}}{\thesubsection}{0.6em}{}
\setlist[itemize]{leftmargin=1.4em,itemsep=0.25em,topsep=0.3em}
\setlist[enumerate]{leftmargin=1.6em,itemsep=0.3em,topsep=0.3em}
\renewcommand{\arraystretch}{1.25}

\newcommand{\ok}[1]{\textcolor{acc}{\textbf{#1}}}
\newcommand{\hold}[1]{\textcolor{warn}{\textbf{#1}}}
\newcommand{\bad}[1]{\textcolor{dead}{\textbf{#1}}}

\title{\vspace{-1.6cm}\textbf{自动代码审计的天花板}\\[0.3em]
\large 项目计划书}
\author{Zhaohe Dong}
\date{2026 年 9 月 17 日}

\begin{document}
\maketitle
\vspace{-0.8cm}

\begin{center}
\begin{minipage}{0.92\textwidth}
\small\itshape
本文件是内部工作计划，不是对外材料。其中对现状的描述包含全部已知的负面结果与已撤回的主张；
凡未通过独立跨厂商评审的数字，均标注其状态，不作为论据使用。
\end{minipage}
\end{center}
\vspace{0.5cm}

\noindent\textbf{这份文件怎么读。}第 1 节讲\emph{初衷}——这条研究线从一个产品开始，
而不是从一个研究问题开始，它最初相信的东西被自己的测量推翻过两次。
第 2 节把研究问题说清楚，第 3 节是装置，第 4 节是术语表与研究编年
（\textbf{不熟悉本项目的读者应先读第 4.1 节的十个词}）。
第 5 节是当前证据状态，包含 2026-09-15 五项评审全部被拒的完整记录。
第 6 节是已识别的五条弱点，第 7、8 节是针对它们的两阶段计划，
第 9 节是论文叙事及其必须遵守的约束，第 10、11 节是交付物与风险。

\vspace{0.3cm}
\noindent\textbf{一句话概括现状}：装置可用、纪律有效、四条主张成立但各自带着严重限定；
最有意思的那条发现证据最弱；上一轮送审的五项研究\emph{全部}被驳回，其中两项是被前一轮的
\emph{修复本身}驳回的。\textbf{这不是严谨的证据，是尚未收敛的证据}——本计划的第一阶段
就是为了让它收敛。
\vspace{0.4cm}

\section{起源：一个产品，和它被自己推翻的前提}

本项目不是从一个研究问题开始的，是从一个产品开始的。这一节说明它的初衷、
它最初相信什么、以及那个信念怎么被自己的测量推翻——\textbf{这是整条研究线的真正起点，
后面每一项研究都能追溯到这里。}

\subsection{CrossAudit 是什么}

一句话定位（写于 2026-08-12 的 \texttt{docs/PRODUCT\_VISION.md}）：

\begin{quote}
\textbf{一个自带独立质量监督体系的 Codex}——用户只描述目标，系统自主完成工作、验证、修订和交付；
只有真正需要人类判断时才打断用户。
\end{quote}

机制是双源审计。一个模型（\emph{生成器}）把文件写进 Git 仓库；确定性检查层（DCL）先跑；
然后\textbf{另一家厂商}的模型（\emph{审计者}）读同一个 commit 并给出 findings，每条标明层级——
是被检查验证的，还是模型提出的。结论只有三种：通过（绑定一张覆盖裁决、规则、commit 与证据集的回执）、
需要修改（生成器带约束改一轮）、需要你（问一个问题并等待）。

\textbf{关键设计：厂商分离由框架强制，不是在提示词里请求的。}这一点是整个产品的卖点，
也正是下面要被检验的那个前提。

\begin{figure}[htbp]
\centering
\begin{tikzpicture}[node distance=7mm]
  \node[box] (goal) {用户描述目标};
  \node[dark, right=9mm of goal] (gen) {生成器\\[-2pt]\footnotesize 厂商 A};
  \node[box, right=9mm of gen] (git) {写入 Git\\[-2pt]\footnotesize commit};
  \node[acc, right=9mm of git] (dcl) {DCL\\[-2pt]\footnotesize 确定性检查};
  \node[dark, right=9mm of dcl] (aud) {审计者\\[-2pt]\footnotesize 厂商 B};

  \draw[fl] (goal) -- (gen);
  \draw[fl] (gen) -- (git);
  \draw[fl] (git) -- (dcl);
  \draw[fl] (dcl) -- (aud);

  \node[acc, below=13mm of git] (pass) {通过\\[-2pt]\footnotesize 绑定回执};
  \node[warnb, below=13mm of dcl] (rev) {需要修改\\[-2pt]\footnotesize 有约束地改一轮};
  \node[box, below=13mm of aud] (ask) {需要你\\[-2pt]\footnotesize 问一个问题};

  \draw[fl] (aud) -- (ask);
  \draw[fl] (aud) -- (rev);
  \draw[fl] (aud.south west) -- (pass.north east);
  \draw[fl] (rev) to[out=180,in=270] (gen.south);

  \node[lbl, below=2mm of pass, align=center]
    {回执绑定：裁决 $\cdot$ 规则 $\cdot$ commit $\cdot$ 证据集};

  \begin{scope}[on background layer]
    \node[draw=dead, dashed, thick, rounded corners=3pt,
          fit=(gen)(aud), inner sep=9pt] (sep) {};
  \end{scope}
  \node[lbl, text=dead, above=1mm of sep.north] {厂商分离由框架强制，不由提示词请求};
\end{tikzpicture}
\caption{CrossAudit 的产品回路。生成器与审计者来自\emph{不同厂商}，
DCL 在审计者之前运行，三种裁决各有出口，通过时绑定一张可复核的回执。}
\end{figure}

成功标准同样写在愿景文件里：用户像用 Codex 一样自然地工作，却能在需要时\emph{证明}
结果由谁生成、检查了什么、为何通过，以及交付物是否与审计证据完全一致。

\subsection{最初的信念，和推翻它的测量}

厂商分离的理由建立在一个直觉上：\textbf{模型看不见自己的缺陷}。自己审自己，等于让考生批自己的卷子。

\textbf{这个直觉是错的，而且是被本项目自己的测量推翻的。}在一项 30 实例的研究里，
自审模型\emph{命名出了更多}缺陷——31.7\% 对陌生审计者的 19.8\%。它不是看不见。

真正的差别在于\textbf{它原谅它们}：自审把 64 条 findings 中的 61 条归为"建议"，
只拦下 3.3\% 的增量；而跨厂商审计者把 48 条 findings \emph{全部}归为阻断，拦下 80.0\%。

\begin{quote}
\textbf{分离买到的是严重度，不是视力。}陌生人是那个不会把你的工作挥手放行的审阅者。
\end{quote}

产品 README 现在就是这么写的。\textbf{这是本项目的第一个范式}：一个卖点建立在一个直觉上，
认真测量之后，卖点还在，但\emph{理由换了}。而"为什么有效"和"是否有效"是两个不同的问题，
只有前者能指导下一步做什么。

\subsection{第二次失败：报告没错，错在引用报告的那一层}

第二个教训来自 \texttt{benchmarks/CORRECTIONS.md}——一份专门记录"哪些数字被撤回或重述"的文件。
它存在的理由写在开头：这些结果打算拿去发表，而一个正在决定要不要信任这份记录的读者，
\textbf{有权在自己发现之前先看到记录出过什么错}。

它的诊断很具体：

\begin{quote}
提交的报告本身是谨慎的——它们给出置信区间、写明 $n$，其中几份还纠正了同一条线上更早的报告。
\textbf{失败发生在报告的下游}，在结果被总结成头条主张的那一层：计划文档、决策记录、
以及工作对话。在那一层，\textbf{区间被丢掉了}，一个"一中之一"的测量以"100\%"的形式流传出去。
有一个数字不只是被剥掉了区间，而是\textbf{被彻底编造出来的}。
\end{quote}

由此产生了本项目最重要的一条制度：\textbf{规则必须绑定在数字被\emph{引用}的地方，
而不只是它被\emph{计算}的地方。}\texttt{EXPERIMENT\_RECORD.md} 随后把这条写成了
每项研究必须满足的落盘标准——写在运行之前的主张、足以重建运行的溯源、
每实例每臂一行的原始记录，而不只是汇总。

\textbf{这两次失败共同决定了后面一切的形态。}第一次教会我们：直觉需要被测量，
哪怕测量结果支持结论、理由也可能是错的。第二次教会我们：报告写对不够，
\emph{引用}报告的那一层才是泄漏发生的地方。

\subsection{从产品问题到研究问题}

两次修正之后，剩下一个产品回答不了的问题：

\begin{quote}
在产审计器在缺陷上的八读并集召回是 \textbf{30.0\%}，代价是正确代码上 \textbf{16.0\%} 的误报。
\textbf{为什么不是 60\%？}继续投入——换更强的模型、改规则手册、多读几遍——能把它推到哪里？
\end{quote}

这个问题不能靠继续调产品来回答，因为分不清"模型不够强"和"题目本身没定义对错"。
\textbf{于是有了这条研究线。}

\section{核心问题与定位}

一个自动审计者读一段不是它写的代码，决定要不要拦下来。它能拦住多少，天花板由什么决定？

这个问题的难点不在提问，在测量。通行做法是让模型评价模型，于是\textbf{“审计者没看见”与
“对错本身没有定义”混在一起}。本项目的全部设计都服务于把这两者分开：让解释器执行的隐藏测试
充当唯一判准，然后一次只操纵一个变量——谁在读、规则手册怎么写、读几次、题目有没有把话说清楚。

\textbf{定位需要收窄。}项目早期的叙事是“AI 审计存在一个天花板”。两轮外部评审都指出这个叙事
撑不住，本计划采纳该判断。可辩护的定位是：

\begin{quote}
\textbf{在给定任务、证据、审计协议与预算下，重复的 LLM 审查能增加多少\emph{有效诊断}？
这些收益付出了多少误报代价，又在多大程度上受任务规格完整性的限制？}
\end{quote}

\section{已建成的装置}

\subsection{判准与分层}

基底为 EvalPlus（HumanEval+/MBPP+），中位题目 41 词规格、6 行参考代码。生成器
\texttt{claude-haiku-4-5} 看规格与\emph{可见}测试写代码；\emph{隐藏}测试只喂给解释器，
审计者永远看不到。分层完全由执行结果决定：

\begin{center}
\begin{tabular}{@{}lll@{}}
\toprule
\textbf{层} & \textbf{条件} & \textbf{规模} \\
\midrule
P（有缺陷） & 过可见测试、挂隐藏测试 & 110 个实例 / 56 个题目 \\
C（“正确”） & 两套测试均通过 & 150 个抽样实例 \\
F（丢弃） & 未过可见测试 & 不进入分析 \\
\bottomrule
\end{tabular}
\end{center}

\begin{figure}[htbp]
\centering
\begin{tikzpicture}[node distance=6mm]
  \node[box] (spec) {规格（散文）\\[-2pt]\footnotesize 中位 41 词};
  \node[box, below=5mm of spec] (vis) {可见测试\\[-2pt]\footnotesize base suite};
  \node[box, right=42mm of spec] (hid) {隐藏测试\\[-2pt]\footnotesize plus suite};

  \node[dark, below=5mm of vis] (gen) {生成器};
  \node[box, below=5mm of gen] (cand) {候选代码};
  \node[acc, below=5mm of cand] (exec) {解释器执行两套测试};

  \draw[fl] (spec) -- (vis);
  \draw[fl] (vis) -- (gen);
  \draw[fl] (gen) -- (cand);
  \draw[fl] (cand) -- (exec);
  \draw[fl] (hid.south) -- ++(0,-30mm) -| (exec.east);

  \node[box, below=6mm of exec, minimum width=62mm] (str)
    {P 有缺陷（110）\quad$\vert$\quad C “正确”（150）\quad$\vert$\quad F 丢弃};
  \draw[fl] (exec) -- (str);

  \node[dark, right=42mm of gen] (aud) {审计者 $\times\,K$\\[-2pt]\footnotesize holistic, $K=8$};
  \node[box, below=16mm of aud] (out) {findings + 严重度\\[-2pt]\footnotesize $\ge 1$ BLOCKER $\Rightarrow$ 拦下};
  \draw[fl] (aud) -- (out);
  \draw[fl] (spec.east) -- ++(12mm,0) |- (aud.west);
  \draw[fl] (cand.east) -- ++(8mm,0) |- (aud.south west);

  \draw[cut] ($(hid.south)+(-9mm,-4mm)$) -- ($(hid.south)+(9mm,-4mm)$);
  \node[lbl, text=dead, right=1mm of hid, yshift=-7mm, align=left]
    {审计者\\永远看不到};

  \node[acc, below=6mm of out, minimum width=44mm] (est) {union recall at $K$};
  \draw[fl] (out) -- (est);
  \draw[fl] (str.east) -- (est.west);
\end{tikzpicture}
\caption{实验装置。审计者读到规格、可见测试与候选代码；\textbf{隐藏测试只喂给解释器}。
这道切断是整个装置成立的前提——它让“没抓到”成为可测量的事件，而不是一个判断。}
\end{figure}

“正确”打引号是认真的：它仅表示“通过隐藏测试”，是解释器给的标签，不是符合规格的保证。
相应地，“误报”指对通过测试的代码提出标记，而非已证明的审计者错误。

\subsection{测量量}

\textbf{union recall at $K$}：$K$ 次独立阅读中至少一次拦下的缺陷实例比例，对全部
$\binom{8}{K}$ 个子集\emph{精确平均}而非抽样。配套量为并集误报率。
方法章立有硬规矩：\textbf{报召回必须同时报它花掉的误报率}——因为全部拦下可得 100\% 召回。

观测单位是实例，重抽单位是题目（同题目的实例不独立）；区间为 10{,}000 次题目级聚类自助法，
Wilson 区间并列打印并标注“过窄”。$K$ 次阅读是重复测量，不是 $K$ 个独立观测。

\subsection{证据闸门}

\begin{enumerate}[label=\textbf{\arabic*.}]
\item \textbf{预注册}：总体、假设、主要结局、区间、kill 条件、预算、停止规则，
      全部在\emph{第一次模型调用之前}写死。
\item \textbf{编号修正案}：任何偏离记为修正案，在其生效的那一步\emph{之前}提交，
      并记录当时已知什么。
\item \textbf{跨厂商评审}：报告交由另一厂商的模型审查，其可读原始记录、重跑分析、
      逐句核对，结论须为 “quotable” 方可放行。
\item \textbf{主张清单}（\texttt{CLAIMS.md}）：列出论文被允许陈述的每一句话及其必须携带的
      标签，同时列出\emph{证据不支持、禁止出现}的说法。
\end{enumerate}

\section{术语与研究编年}

\subsection{读本文需要的十个词}

\begin{center}
\begin{tabular}{@{}lp{10.2cm}@{}}
\toprule
\textbf{词} & \textbf{含义} \\
\midrule
生成器 & 写候选代码的模型。全程为 \texttt{claude-haiku-4-5}，冻结 \\
审计者 & 读代码并给出 findings 的模型。它\emph{永远看不到隐藏测试} \\
\texttt{cross} & 跨厂商审计者，\texttt{gpt-5.6-terra}。在产配置 \\
\texttt{self} & 同厂商同规格审计者，\texttt{claude-haiku-4-5}。自己审自己 \\
\texttt{self-strong} & 同厂商更强模型，\texttt{claude-sonnet-4.6} \\
holistic & 一次性读完整个增量的审计形态。在产配置 \\
decomposed & 把审计拆成若干子问题分别提问的形态 \\
DCL & 确定性检查层。不调用模型的规则检查，先于审计者运行 \\
BLOCKER & 最高严重度。\emph{至少一条 BLOCKER} 即判定拦下 \\
回执 & 绑定裁决、规则、commit 与证据集的记录，写进 Git \\
\bottomrule
\end{tabular}
\end{center}

\subsection{研究编年}

本研究线共二十余项研究，下表只列与本计划直接相关的。
\textbf{编号是历史顺序，不是重要性顺序。}

\begin{center}
\begin{tabular}{@{}llp{7.4cm}@{}}
\toprule
\textbf{研究} & \textbf{记录} & \textbf{它问了什么 / 得到什么} \\
\midrule
study 2 & \texttt{RESULTS-2.md} & 拆解审计是否提高召回？是——高 9.1 点，但误报也动了 \\
study 5 & \texttt{RESULTS-5.md} & \textbf{前提检验}：跨厂商审计者命名的缺陷\emph{更少}，
  推翻"看不见自己缺陷"的直觉 \\
study 7 & \texttt{RESULTS-EXPLORE.md} & \textbf{十五种架构一条 ROC}：在产架构就在 Pareto 前沿，
  没有一种在 $\le6.7\%$ 误报下胜过它。\emph{预注册的空结果} \\
study 16 & \texttt{RESULTS-TESTGEN.md} & 会写测试的审计者能否胜过只读的？ \\
study 17 & \texttt{RESULTS-TESTGEN-VAL.md} & 能否在没有标准答案时认出\emph{写错的}测试？
  \emph{未达预注册门槛} \\
ceiling 1 & \texttt{RESULTS-CEILING.md} & \textbf{主研究}：union-of-$K$ 曲线、误报代价、
  规则手册对照。C1/C5/C6 的证据 \\
study 18 & \texttt{RESULTS-CEILING3.md} & 再加两个审计者族。C2 的证据 \\
study 19 & \texttt{RESULTS-CEILING3B.md} & 改了规则手册能否推动第二个厂商的族？ \\
study 20 & \texttt{RESULTS-CEILING4.md} & \textbf{标记是不是抓到？}保留 finding 文本做裁定 \\
study 21 & \texttt{RESULTS-RERATE.md} & \textbf{残差重分类}：先问"规格定了吗"再问别的。C4 的证据 \\
study 22 & \texttt{RESULTS-INJECT.md} & C4 的前瞻性检验。\textbf{已整个撤回} \\
study 23 & \texttt{RESULTS-SUBSTRATE2.md} & 第二个基底。\textbf{整轮作废，重跑中} \\
\bottomrule
\end{tabular}
\end{center}

\textbf{值得注意的是这条线的形状}：它不是一路推进，而是反复自我推翻。study 5 推翻了产品的
原始前提；study 7 得到一个预注册的\emph{空结果}并接受了它；study 21 推翻了更早的残差分类；
study 22 与 23 被评审整个否掉。\textbf{一条只前进不后退的研究线，通常是没有在测量。}

\section{当前证据状态}

\subsection{已通过评审的发现}

\begin{itemize}
\item \ok{C1}\; 重复阅读收益递减但\emph{未饱和}。$K{=}1$ 的 10.7\% 升至 $K{=}8$ 的
      30.0\% [20.0, 40.7]，误报 16.0\% [10.1, 22.3]。第八次仍增益 1.93 点 [1.14, 2.78]，
      高于预注册的 1.0 点平坦化门槛，故\textbf{“饱和”一词已撤下}。三族二十次阅读合计
      48.2\% [36.7, 60.0]。
\item \ok{C2}\; 换模型的影响大于换次数，\textbf{但该对比过不了自身的关}：更强的同厂商模型
      召回 $-26.4$ 点 [$-37.3$, $-15.6$]，然而它误报仅 3.3\%（对 16.0\%）；换宽松严重度规则
      符号\emph{反转}；两条路由采样方式不同（一条发 temperature 0，一条不发），$K{=}1$ 时
      差距仅 $-10.1$。且存在反例：另一跨厂商模型在\emph{一半}阅读次数下两轴皆优。
\item \ok{C3}\; 规则手册加一句话使标记从 10/56 增至 25/56（$+26.8$ 点 [13.6, 40.4]），
      代价为正确代码上 $+12.50$ 点；但在隔离该规则的对比上，
      隐藏测试通过率仅动 $+5.36$ 点 [$-0.89$, 12.07]，$p=0.146$，\textbf{区间跨零}。
\item \ok{C4}\; 110 个缺陷中 57 个被所有阅读漏过；重新分类后两位评分者一致认为
      \textbf{44/57（77.2\% [62.1, 91.1]）是规格未确定隐藏测试所问之事}。
      外部佐证：他人独立点名的 11 道题目中，我方冻结标签吻合 9 道。
\end{itemize}

\subsection{2026 年 9 月 15 日：五项评审，五项全拒}

\begin{center}
\begin{tabular}{@{}lll@{}}
\toprule
\textbf{研究} & \textbf{状态} & \textbf{缘由} \\
\midrule
study 17 & \hold{已修} & 第 8 轮的修复立了新过强主张：声称测试证明了一条从未执行的重建链 \\
study 19 & \hold{已修} & 改名只做一半；且把工具准入说成“更窄”，实际更宽 \\
study 20 & \hold{已修} & 把八读\emph{标记覆盖}写成已确立的更高\emph{天花板} \\
study 22 & \bad{已撤回} & 注册要复原失败输入的过滤器，实现只查字符串非空 \\
study 23 & \bad{作废·重跑中} & 给模型看的可见测试缺 class 头，300 道题全部语法损坏 \\
\bottomrule
\end{tabular}
\end{center}

\textbf{两项是被上一轮修复本身拒掉的}，且偏差均指向“让工作显得更好”的方向。这不是严谨的证据，
是尚未收敛的证据。所幸“未过评审的数字不进正文”这条规矩，使五连拒\emph{未逼出任何已认定主张的
撤回}——论文的四个结果小节依赖的三项研究都是先过评审才进来的。

\section{已识别的最大弱点}

按严重程度排列，本计划的下一阶段直接针对前三条。

\begin{enumerate}[label=\textbf{W\arabic*.}]
\item \textbf{C4 的评分者不独立。}这是最有意思的发现，也是防护最薄的一条。两位评分者是
      一个模型与\emph{作者本人}——后者写了评分标准、知道先前标签、知道该主张的分量。
      且“规格是否确定”本质上是判断，解释器管不了。为其设计的前瞻性检验（study 22）已整个撤回，
      因此 C4 目前\textbf{只有事后证据}。
\item \textbf{对照组已做，但从未进入论文。}残差重分类之外，两位评分者也标注了 53 个
      \emph{已抓到}的缺陷（\texttt{L1-flagged.csv} / \texttt{L2-flagged.csv}，经核查为独立评分
      而非复制）。共识标签下：\textbf{漏检组 46/68 = 67.6\% 规格未确定，已抓到组
      24/53 = 45.3\%，差 $+22.4$ 点}；两组评分者一致率分别为 88.2\% 与 96.2\%。
      这是 C4 现在缺的那个对照，\textbf{它支持 C4}——漏检确实更集中在规格未确定的实例上，
      但并非全部。该对比至今未写进任何报告，也未过评审，须补上区间与标签后送审。
\item \textbf{关键实验未做。}严重度阈值扫描——把不同审计者拉到\emph{同一误报率}再比——
      决定 C2 与 C5 能否成立，而所需 findings 均已归档，\textbf{只改评分规则、几乎不花钱}。
\item \textbf{单基底单生成器。}第二基底作废、注入研究撤回后，论文只剩一个基底、一个生成器支撑，
      且该基准记忆污染问题被广泛批评。
\item \textbf{区间覆盖未验证。}主区间为 56 个题目簇上的自助法，其在本聚类设计下的覆盖率
      未经模拟验证；本项目此前已有两个区间构造被评审测出覆盖率为 0.416 与 0.075。
\end{enumerate}

\section{下一阶段工作}

\subsection{P1 —— 人类评分者与对照组（针对 W1、W2）}

\textbf{这是性价比最高的一项，且不需要任何模型调用。}

设计：把 110 个缺陷实例\emph{全部}交给评分者，\textbf{对审计结果盲化并打乱顺序}，
只问一个问题——“规格的散文是否确定了隐藏测试所期望的值？”

\begin{itemize}
\item \textbf{主要对照已有初值，需独立复核}：已抓到与漏检两组“规格未确定”的比例差为
      $+22.4$ 点（45.3\% 对 67.6\%）。但该数字的两位评分者\emph{仍是}一个模型与作者本人，
      因此它证实的是 C4 的方向，不是 C4 的独立性。P1 的任务是让一位项目外的人类在盲化条件下
      重做这个对比。
\item \textbf{附带产出}：作者对同 57 个实例的\emph{重测信度}。该评分标准已经因提问顺序
      改变而翻转过一次，其稳定性本身是值得记录的量。
\end{itemize}

\textbf{独立性限制必须写明}：作者曾作为 L1 评过那 57 个漏检实例，故该部分是\emph{重测}而非
独立评审；53 个已抓到实例是新的。真正的独立评审仍需一位项目外的人类，哪怕只标 30 项。

\subsection{P2 —— 严重度阈值扫描（针对 W3）}

在单一基底上对每个审计者族生成召回--误报曲线，所有跨族比较在\emph{匹配的误报率}上进行。
现有两个阈值点（BLOCKER、任意 finding）已在手，补第三、第四个阈值只需对已缓存的 findings
重新评分。\textbf{这是决定 C2/C5 生死的实验，却几乎免费。}

\subsection{P3 —— 规格澄清实验（针对 W1 的因果性）}

保持代码与可见测试\emph{完全不变}，只补上规格中从未交代的行为规则（同分如何处理、
空输入返回什么），并配一个长度相当、但不解决任何歧义的\textbf{安慰剂改写}作为对照。

\begin{center}
\begin{tabular}{@{}llll@{}}
\toprule
\textbf{条件} & \textbf{代码} & \textbf{可见测试} & \textbf{规格} \\
\midrule
原始 & 不变 & 不变 & 原始描述 \\
澄清 & 不变 & 不变 & 补足缺失的行为规则 \\
安慰剂 & 不变 & 不变 & 增加等长但不解决歧义的文字 \\
\bottomrule
\end{tabular}
\end{center}

若正确诊断率在澄清条件下上升而在安慰剂下不升，规格的不完整就在起\emph{因果}作用；若不升，
则 C4 只是“难实例的两个属性相关”。\textbf{该设计操纵规格，而重分类与注入总体都只是在规格上
做条件——这是它们做不到的。}

\subsection{P4 —— 收尾 study 23}

重跑正在进行中。self 族 8 臂已全部完成；cross 族因供应商限流中断约 23 小时后已恢复，
正在补齐。预算封顶 \$45（预注册 Amendment 2），当前约 \$30。

\section{第二阶段：审计形态与 AI4S}

第一阶段（P1--P4）是把现有主张钉牢。第二阶段是往外走，分两步：\textbf{先广泛调研当前已有的
审计架构与方法，再结合 AI4S 数据集做能落地的分析}。这个顺序有两个理由：外部评审明确指出本项目
文献 engagement 不足（现有引用仅 11 条），调研能直接补上；而先看清别人做过什么，才能避免把预算
花在已被回答的问题上——本项目自己的 study 7 就是例子。

\subsection{P0' —— 先做调研，再做实验}

\textbf{已知的自身空白}：study 7 已经用十五种架构（holistic/decomposed $\times$
cross/self/cheap-cross $\times$ union/majority/多 draw）在 $\le 6.7\%$ 误报约束下做过网格搜索，
结论是\emph{没有一种胜过在产的} \texttt{holistic-cross}，它就在 Pareto 前沿上
（\texttt{tri\_union} 召回 41.8\% 但误报 33.8\%，换不回来）。
\textbf{因此不应再做架构排列组合，边际收益已被自己证明很低。}

调研应集中在那十五种\emph{没有覆盖}的轴上，并在每条上确认是否已有他人工作：

\begin{center}
\begin{tabular}{@{}lp{7.6cm}@{}}
\toprule
\textbf{未覆盖的轴} & \textbf{为什么值得} \\
\midrule
工具使用（可执行代码） & 十五种全是只读式。\textbf{我们测到的可能是“阅读”的天花板，不是“审计”的天花板} \\
顺序 / 自适应 & 十五种全是并行合成，没有接力：审计者 2 看得到审计者 1 的 findings \\
性质导向（写不变量而非 findings） & 输出可执行 $\Rightarrow$ 判准不再是判断 \\
辩论 / 对抗 & 有文献基础，但成本高，需先看已有结论 \\
检索增强 & 给审计者看相似的已知缺陷 \\
\bottomrule
\end{tabular}
\end{center}

交付物是一份带引用的对照表：\emph{每条轴上别人做到哪、用什么判准、我们能加什么}。
这份表同时是论文相关工作章的材料。

\subsection{P5 —— 只读式 vs 工具辅助审计}

\textbf{架构侧唯一的大空白，也是二选一的干净对照。}给审计者一个解释器：它可以自己写测试、
自己跑、自己看崩在哪；但\emph{不给}隐藏测试、不给标准答案。其余一切（基底、实例、规格、
可见测试、$K$、严重度规则）与 ceiling 1 冻结的配置相同。

回答的问题是：\textbf{天花板是模型的性质，还是形态的性质？}若工具辅助显著突破 30.0\%，
则本项目全部已认定主张的适用范围都必须收窄到\emph{只读式代码审查}——这正是论文当前声明的范围，
届时该声明将从“保守的自我设限”变成“已被测量的边界”。

预注册须事先写死：工具预算（调用次数、墙钟）、可执行环境、以及\textbf{“审计者自己写的测试
通过与否”不得进入判准}——判准仍然只有隐藏测试。

\subsection{P6 —— 科学代码中“规格未确定”的比例}

\textbf{不在 AI4S 基底上重做天花板研究，只搬那个诊断动作。}假设是：

\begin{quote}
科学代码的规格天然更不完整——\textbf{容差取多少算收敛、单位、随机种子、数值稳定性}
极少写进题面——因此“规格未确定”的比例应\emph{显著高于}玩具 Python 的 67.6\%。
\end{quote}

若成立，这是一条直接给 AI4S 领域的结论：\textbf{继续堆模型的边际收益有上限，而上限由任务定义
的质量决定。}它不依赖本项目的模型、工具链或基准。

\textbf{候选基底与两个已知障碍。}\emph{SciCode}（NeurIPS 2024 D\&B）覆盖物理、数学、材料、
生物、化学五个领域，80 个主问题拆为 338 个子问题，带科学家标注的金标准解与测试用例。
但直接套用本项目的装置有两个硬问题，\textbf{必须在预注册里先解决}：

\begin{enumerate}[label=\textbf{障碍 \arabic*：}]
\item \textbf{它没有“可见/隐藏”测试划分。}substrate 2 正是因为\emph{自己发明}这个划分而整轮
      作废——切出来的可见测试 300 道题全部无法解析。\textbf{同一个坑不得再踩。}
\item \textbf{求解率过低。}最好的模型在最真实设定下仅解出 7.7\%，绝大多数候选会落进丢弃层，
      P/C 两层都凑不出样本。
\end{enumerate}

\textbf{因此 P6 的设计绕开生成与审计}：不生成候选、不跑审计，只取现成的失败案例与题面，
交由评分者判断“规格是否确定了失败处应返回什么”。成本接近零，且不触碰上述两个障碍。

\subsection{P7 —— 不变量审计（科学代码独有的机会）}

科学计算有一样玩具 Python 没有的东西：\textbf{物理不变量}——能量守恒、动量守恒、对称性、
极限情形、量纲一致性。它们是判准，\emph{且不依赖任何测试套件}。

设计：不问审计者“这段代码有缺陷吗”，而问\textbf{“它违反了哪条不变量”}，
并要求它给出\emph{可执行的检查}。该检查随后由解释器运行。

\textbf{这等于绕开了规格确定性问题}——不再需要题面把话说死，因为物理定律替它说了。
同时它也是 P0' 表中“性质导向”那一行的具体实现：审计者的输出从判断变成可执行断言。

\textbf{风险}：不变量可能过于宽松（守恒成立但结果仍错），也可能过于严格（数值误差触发假阳）。
\textbf{因此 P7 须先做小规模可行性验证}——在 20--30 个已知缺陷上确认不变量检查的区分度，
再决定是否扩大。不通过则记录并停止，不追加预算。

\section{预算与时间}

\begin{center}
\begin{tabular}{@{}lllp{5.2cm}@{}}
\toprule
\textbf{项} & \textbf{模型开销} & \textbf{人力} & \textbf{说明} \\
\midrule
P1 人类评分与对照 & \$0 & 约 4--6 小时 & 无模型调用；盲化数据集由脚本生成 \\
P2 阈值扫描 & 约 \$0 & 约 1 天 & 仅对已归档 findings 重新评分 \\
P3 规格澄清实验 & 约 \$40--60 & 约 3--5 天 & 需新生成与新审计；应单独预注册 \\
P4 study 23 收尾 & 余 \$15 以内 & 已自动化 & 守护进程运行中 \\
\midrule
\multicolumn{4}{@{}l@{}}{\textit{第二阶段}} \\
P0' 调研 & \$0 & 约 3--5 天 & 五条未覆盖轴的带引用对照表；同时是相关工作章的材料 \\
P5 只读 vs 工具辅助 & 约 \$60--100 & 约 1 周 & 工具调用使单次审计成本上升，需单独预注册预算上限 \\
P6 科学代码规格比例 & \$0 & 约 1 周 & 不生成、不审计，只评分；绕开 SciCode 的两个障碍 \\
P7 不变量审计 & 约 \$15（可行性） & 约 1 周 & 先在 20--30 例上验区分度，不通过即停 \\
\bottomrule
\end{tabular}
\end{center}

\textbf{排序理由}。P1 与 P2 几乎不花钱却直接打在最弱的两环上，应先做；P3 是唯一能把相关性
变成因果主张的实验，但需要预算与新的预注册，应在 P1 结果出来之后再启动——若 P1 的对照显示
漏检组与已抓到组比例相近，P3 的前提就不成立，应当改题而非照做。

第二阶段的顺序是\textbf{先调研、再落地}：P0' 不花钱，却能防止把预算花在已被回答的问题上——
本项目自己的 study 7 就是这样一个例子，十五种架构的网格搜索得到的是一个空结果。
P6 同样零成本且不依赖第一阶段的任何结论，可与 P0' 并行。
\textbf{P5 与 P7 都要花钱，都应等 P0' 的对照表出来之后再启动}：若调研显示某条轴已被他人
用更强的判准回答过，该项应改为复现或引用，而不是重做。

\begin{figure}[htbp]
\centering
\begin{tikzpicture}[node distance=6mm]
  \node[acc] (p1) {\textbf{P1} 盲化人工评分\\[-2pt]\footnotesize \$0 · 打在 W1/W2};
  \node[acc, right=8mm of p1] (p2) {\textbf{P2} 阈值扫描\\[-2pt]\footnotesize \$0 · 定 C2/C5 生死};
  \node[warnb, right=8mm of p2] (p3) {\textbf{P3} 规格澄清实验\\[-2pt]\footnotesize \$40--60 · 唯一因果设计};
  \draw[fl] (p1) -- (p3);
  \node[lbl, above=0.5mm of p3, align=center] {若 P1 显示两组接近\\则前提不成立，应改题};

  \node[acc, below=15mm of p1] (p0) {\textbf{P0$'$} 调研\\[-2pt]\footnotesize \$0 · 五条未覆盖轴};
  \node[acc, below=15mm of p2] (p6) {\textbf{P6} 科学代码规格比例\\[-2pt]\footnotesize \$0 · 可并行};
  \node[warnb, below=15mm of p3] (p5) {\textbf{P5} 只读 vs 工具辅助\\[-2pt]\footnotesize \$60--100};
  \node[warnb, below=8mm of p5] (p7) {\textbf{P7} 不变量审计\\[-2pt]\footnotesize \$15 可行性};

  \draw[fl] (p0) -- (p5);
  \draw[fl] (p0) -- (p7.west);

  \node[box, left=8mm of p1, rotate=90, anchor=north, minimum width=22mm] (a1) {第一阶段};
  \node[box, left=8mm of p0, rotate=90, anchor=north, minimum width=22mm] (a2) {第二阶段};

  \node[lbl, below=1.5mm of p0, align=center, text width=34mm]
    {不花钱，却能防止把预算\\花在已被回答的问题上};
\end{tikzpicture}
\caption{两阶段的依赖与排序。绿框零成本、应先做；橙框要花钱、都有前置条件。
箭头表示\emph{前置依赖}而非时间顺序——P6 与 P0$'$ 可与第一阶段并行。}
\end{figure}

\textbf{一条跨阶段的硬规矩}：任何新基底在第一次模型调用之前，必须先跑一次
\textbf{“模型将看到的文本能否解析”}的闸门。这是 substrate 2 用一整轮作废换来的教训——
当时可见测试文件 300 道题全部无法解析，而审计者与生成器都读了它。

\section{论文叙事}

拟采用三幕结构，与项目的真实发生顺序一致：

\begin{enumerate}[label=\textbf{第 \arabic* 幕}]
\item \textbf{从产品出发。}我们造了一个跨厂商代码审计器并把它发出去了。认真测量之后，
      它在缺陷上的八读并集召回是 30.0\%，代价是正确代码上 16.0\% 的误报；
      而且\emph{标记不等于抓到}——一次阅读里它对 110 个缺陷提出 9 条 finding，
      真正说对失败行为的只有 5 到 7 条。
\item \textbf{追问天花板。}于是把“审计者能抓到多少”拆开：谁在读、规则手册怎么写、读几次、
      题目有没有把话说清楚。四条发现在此，包括最反直觉的那条——
      漏掉的东西里 67.6\% 是\emph{规格根本没规定隐藏测试在问什么}。
\item \textbf{到 AI4S 去检验。}把这套诊断搬到科学代码上：那里的规格更不完整（容差、单位、
      收敛判据、数值稳定性极少写进题面），因此“规格未确定”的比例应当更高。
      若成立，这是一条对 AI4S 方向直接有用的结论。
\end{enumerate}

\subsection{为什么这个叙事比原来的好}

原叙事是“AI 审计存在一个天花板”。两轮外部评审都指出它撑不住：三项发现在文献里已有先例，
而“天花板”一词被自己的预注册门槛否掉。三幕结构解决了三个问题——它给测量提供了\emph{真实的动机}
（产品没达到预期，所以去查为什么），它让负面结果成为\emph{论据而非尴尬}
（一篇讲“我们试过、这是学到的”的论文可以自然吸收撤回的研究），
它把射程从玩具 Python 延伸到一个有人关心的领域。

\subsection{三条必须先钉死的约束}

\textbf{一、不要说“失败的 CrossAudit”。}记录不支持这个词。study 7 用十五种架构做过网格搜索，
在产架构\emph{就在 Pareto 前沿上}，没有一种在误报约束下胜过它。所以它不是做砸了的产品，
而是\textbf{一个按设计工作、但天花板比预期低的产品，并且低的原因有一大半不在它身上}。
“失败”是一个比“已测量的天花板”更弱也更不准确的说法。诚实的措辞是：
\emph{我们发出了一个审计器，认真测了它，发现天花板不在我们以为的地方。}

\textbf{二、叙事不得重排证据。}这是本项目最大的特有风险。整套纪律建立在“假设先于数据”之上，
而一个“从失败出发、发现了 X”的故事，天然诱使人把一个更混乱的真实过程整理成一条干净的因果链。
论文必须继续逐条标注哪些是预注册主分析、哪些是探索性、哪些是看过结果之后的修订。
\textbf{\texttt{CLAIMS.md} 的存在正是为了防止这件事，叙事不能凌驾于它。}

\textbf{三、第三幕目前还不存在。}P6 与 P7 尚未运行。因此当下有两个选择：
先跑 P6 再写第三幕，或者把第三幕写成\emph{计划}而非\emph{结果}并明确标注。
\textbf{不可以把尚未做的实验写成已有的发现}，哪怕只是叙事上的暗示。
建议取前者——P6 零成本，是三项里最便宜的一项。

\section{交付物}

\begin{itemize}
\item \textbf{近期}：P1 的盲化评分数据集与结果表；P2 的召回--误报曲线；study 23 的重跑报告
      （按 Amendment 2 与作废轮并列，不附 $p$ 值或区间）。
\item \textbf{论文}：当前 7 页，ICML 风格预印本。\textbf{硬性阻断}——四项在评审中的研究
      全部通过、或其数字全部移除之前，不上 arXiv、不投任何会议。
\item \textbf{第二阶段}：P0' 的带引用对照表（五条未覆盖轴）；P5 的只读 vs 工具辅助对照，
      它决定现有主张的适用范围是否必须收窄到只读式审查；P6 的科学代码规格未确定比例，
      这是一条可直接交付给 AI4S 方向的结论。
\item \textbf{可复用产出}：本项目认为价值高于论文本身的，是那个\emph{诊断动作}——
      把失败案例重新分类，区分“模型的局限”与“需求的沉默”，并报出比例。
      该动作不依赖本项目的模型、基准或工具链。
\end{itemize}

\section{风险}

\begin{itemize}
\item \textbf{P1 的对照可能证伪 C4。}这是设计它的理由，不是它的风险。若发生，
      应如实报告并改写论文主张，而非换一种切法直到结果顺眼。
\item \textbf{评分标准不稳定。}已知它会随提问顺序翻转。P1 的重测信度正是为量化这一点。
\item \textbf{供应商可用性。}study 23 已因限流停摆约 23 小时。所有长跑须在耐久目录下
      由守护进程承载，并设中途停机线。
\item \textbf{循环论证。}研究对象是 LLM，评审者、评分者之一亦是 LLM。缓解手段是：
      判准为解释器而非模型；预注册；跨厂商评审。\textbf{但在判准为“判断”的地方——恰是
      C4——该缓解不成立}，这正是 P1 与 P3 存在的理由。
\end{itemize}

\end{document}
